\chapter{Fuchs--Painlev\'e Algebraic Differential Equations}
\label{ch:algebraic-ode}

\section{Scope and foundation}

This subproject uses the project's rational interval algorithms, computable
complex values, and explicit finite-difference derivative certificates.
There is no Mathlib real-number dependency. The first checked milestone
concerns the Euler--Fuchs equation and Painlev\'e I and II; the general
Fuchs--Painlev\'e classification remains a research goal.

\begin{definition}[Polynomial differential residuals]
\label{def:algebraic-ode-residual}
A finite expression in $x,y,y',y''$, with rational constants, addition,
negation and multiplication, can be evaluated over rational numbers,
rational complex numbers, or raw computable complex algorithms.
An algebraic relation in $x,y$ includes a rational evaluation witness proving
that its polynomial is not identically zero.
\lean{ComputableAnalysis.AlgebraicODE.Expr}
\lean{ComputableAnalysis.AlgebraicODE.AlgebraicRelation}
\leanok
\end{definition}

\begin{theorem}[Represented residual semantics]
\label{thm:algebraic-ode-runtime}
Evaluating a polynomial expression at valid raw inputs yields a valid raw
output and respects overlap equivalence. At exact rational inputs its
runtime is the exact rational evaluation at every stage.
This is a residual theorem: derivative certificates remain separate.
\lean{ComputableAnalysis.AlgebraicODE.Expr.evalRaw_valid}
\lean{ComputableAnalysis.AlgebraicODE.Expr.evalRaw_equiv}
\lean{ComputableAnalysis.AlgebraicODE.Expr.evalRaw_ofRat}
\leanok
\end{theorem}

\section{Euler--Fuchs polynomial solutions}

\begin{theorem}[Coefficient classification]
\label{thm:fuchs-polynomial-classification}
For $xy'=r y$ and a finite rational polynomial $p=\sum_k c_kx^k$,
the coefficient equations are $(k-r)c_k=0$ for every $k$.
They hold exactly when each nonzero coefficient has $k=r$, and imply
$x p'(x)=r p(x)$ under the existing executable polynomial derivative.
In particular, a nonzero polynomial solution forces $r$ to be a
nonnegative integer.
\lean{ComputableAnalysis.AlgebraicODE.Fuchs.polynomial_classification}
\lean{ComputableAnalysis.AlgebraicODE.Fuchs.polynomial_equation}
\lean{ComputableAnalysis.AlgebraicODE.Fuchs.exponent_of_nonzero_polynomial}
\leanok
\end{theorem}

For a radical relation $y^n=x^m$, the checked algebraic elimination derives
$nxy'=my$ from the differentiated relation, provided $y\ne0$ and $n\ne0$.
It does not construct a branch or prove implicit differentiation.
\lean{ComputableAnalysis.AlgebraicODE.Fuchs.radical_euler}
\leanok

\section{Painlev\'e equations and algebraic seeds}

We use the DLMF convention
\[
  \mathrm{P_I}: y''=6y^2+x,\qquad
  \mathrm{P_{II}}(\alpha): y''=2y^3+xy+\alpha.
\]
PI has no rational constant solution. For PII, the only rational affine
solution is $y=0$, at $\alpha=0$.
\lean{ComputableAnalysis.AlgebraicODE.Painleve.first_no_constant}
\lean{ComputableAnalysis.AlgebraicODE.Painleve.affine_classification}
\leanok

\begin{theorem}[Simple-pole classification and analytic certificates]
\label{thm:painleve-simple-pole}
On the punctured rational line, $y=c/x$, with
$y'=-c/x^2$ and $y''=2c/x^3$, satisfies PII exactly when
\[
  \alpha=-c,\qquad c\in\{0,1,-1\}.
\]
On every positive rational interval the constructed solution package
contains both two-sided finite-difference derivative certificates, the ODE
identity, and the nonzero algebraic relation $xy-c=0$ for the same evaluator.
The represented-complex residual also vanishes at every nonzero rational
input. Thus these are checked algebraic seeds, not merely assigned jets.
\lean{ComputableAnalysis.AlgebraicODE.Painleve.simple_pole_classification}
\lean{ComputableAnalysis.AlgebraicODE.Painleve.poleSolution}
\lean{ComputableAnalysis.AlgebraicODE.Tests.negativePole_represented}
\leanok
\end{theorem}

The sign symmetry and the Riccati reduction at $\alpha=1/2$ are checked
algebraic identities. The latter explicitly assumes differentiated Riccati
data. Arbitrary-degree polynomial, rational, and algebraic solution
classifications have not yet been proved.

\section{Existing work and next milestones}

A search of public web and GitHub results on 22 September 2026 did not
locate an inspectable formalization of the Fuchs--Painlev\'e algebraic-solution
classification. This is a search result, not a nonexistence claim.

\begin{itemize}
\item \href{https://users-cs.au.dk/spitters/Picard.pdf}
{Makarov and Spitters: the Picard algorithm in Coq} uses constructive
exact reals from CoRN/MathClasses, with
\href{https://github.com/EvgenyMakarov/corn/tree/master/ode}{public ODE source}.
It is the closest foundation-level comparison, although its completion and
function-space approach differs from our finite interval certificates.
\item \href{https://isa-afp.org/entries/Ordinary_Differential_Equations.html}
{Isabelle AFP: Ordinary Differential Equations} provides existence,
uniqueness, flows, and associated verified numerical algorithms. It is a
substantial comparison for analytic ODE specifications, with a different
number foundation.
\item \href{https://cfhp.univ-lille.fr/files/students/Master-2026-Coqodi.pdf}
{Solving Differential Equations in Rocq} is a spring 2026 research proposal
that explicitly includes Painlev\'e I among target examples. It concerns
validated numerical approximation; the proposal is not evidence of a
completed implementation or algebraic classification.
\item \href{https://arxiv.org/abs/2408.15180}
{Mason--Stothers and its corollaries in Lean 4} supplies an adjacent
comparison for polynomial derivative and degree arguments. It is not a
Painlev\'e formalization.
\item \href{https://dlmf.nist.gov/32}{DLMF Chapter 32} supplies the equation
and seed conventions, rather than machine-checked proofs.
\end{itemize}

Next milestones are normalized polynomial/rational differential arithmetic,
certified algebraic branches and Puiseux charts, regular-singular Fuchs
theory, and a separate first-order Fuchs criterion development. Painlev\'e
work will next address pole/degree obstructions, B\"acklund transformations,
and rational-solution classification for PII, before III--VI and their
parameter-dependent algebraic solutions. Analytic continuation, monodromy,
and the general algebraic classification are open.

The new modules add no custom axioms or proof placeholders. Exact
classification proofs use standard Lean logical axioms; raw validity and
derivative packages inherit the foundation's existing native-decision
dependencies, exposed by the subproject's axiom audit.
